\documentclass[journal]{IEEEtran}
\usepackage{cite}
\usepackage{amsmath,amssymb,amsfonts}
\usepackage{algorithmic}
\usepackage{graphicx}
\usepackage{textcomp}
\usepackage{xcolor}
\usepackage{booktabs}
\usepackage{hyperref}

\begin{document}

\title{Autonomous Navigation and Localization for UAVs in GPS-Denied Environments: A PRISMA 2020 Systematic Literature Review}

\author{Abhishek~Raj,~\IEEEmembership{Member,~IEEE}}

\markboth{IEEE Transactions on Robotics,~Vol.~XX, No.~X, June~2026}%
{Raj: Autonomous Navigation for UAVs in GPS-Denied Environments}

\maketitle

\section{GPS-Denied UAV Navigation: A Systematic Literature Review of
Multi-Sensor Fusion, Vision-Based Localisation, and Cooperative Autonomy
(2013--2026)}\label{gps-denied-uav-navigation-a-systematic-literature-review-of-multi-sensor-fusion-vision-based-localisation-and-cooperative-autonomy-20132026}

\textbf{Target Publication:} \emph{IEEE Transactions on Robotics} /
\emph{IEEE Access}\\
\textbf{Corpus Evidence Base:}
\texttt{02\_data\_processed/MASTER\_EVIDENCE.csv} (\(n = 279\)
peer-reviewed studies)\\
\textbf{Quality Framework:} Plan-Do-Check-Act (PDCA) Multi-Tier Quality
Screening
(\texttt{06\_analysis/outputs/quality\_appraisal\_scored.csv})\\
\textbf{Data Analytics Reference:}
\texttt{06\_analysis/outputs/RQ\_DATA\_ANALYTICS.json}

\begin{center}\rule{0.5\linewidth}{0.5pt}\end{center}

\subsection{Abstract}\label{abstract}

The deployment of Unmanned Aerial Vehicles (UAVs) in Global Navigation
Satellite System (GNSS)-denied environments represents one of the most
formidable frontiers in modern robotics. Whether navigating subterranean
mines, inspecting collapsed infrastructure, or operating within
contested airspace, aerial platforms must execute high-frequency control
and complex path planning in the complete absence of global positioning
updates. This Systematic Literature Review (SLR) provides a
comprehensive synthesis of 279 gold-standard, peer-reviewed empirical
studies published between 2013 and 2026. Executed through a rigorous
PRISMA 2020 protocol across IEEE Xplore and Scopus, this review
investigates four fundamental dimensions of GPS-denied flight:
(\textbf{RQ1}) the historical evolution of sensor modalities and fusion
architectures; (\textbf{RQ2}) the empirical performance envelopes across
extreme operational environments, including dense forest canopies and
urban canyons; (\textbf{RQ3}) the validation fidelity of leading
algorithmic paradigms, contrasting the maturity of LiDAR-inertial
systems against the simulation dependency of cooperative swarms; and
(\textbf{RQ4}) the systemic vulnerabilities imposed by Size, Weight,
Power, and Cost (SWaP-C) constraints. Our chronological analysis exposes
a paradigm shift from legacy ultrasonic odometry (2013--2016) to the
contemporary dominance of lightweight solid-state LiDAR and
Ultra-Wideband (UWB) networks (2022--2026). More critically, our quality
appraisal reveals a profound reproducibility crisis: fewer than 2\% of
investigations provide open-source code and verifiable flight datasets.
By aggressively avoiding artificial statistical pooling and relying
exclusively on verbatim empirical metrics, this review dismantles the
theoretical ``Sim-to-Real'' disconnect and defines a concrete,
hardware-grounded roadmap toward certifiably autonomous, closed-loop
aerial robotics.

\textbf{Index Terms}---UAV navigation, GNSS-denied localization,
multi-sensor fusion, visual-inertial odometry, LiDAR SLAM, cooperative
swarms, PRISMA systematic literature review.

\begin{center}\rule{0.5\linewidth}{0.5pt}\end{center}

\subsection{1. Introduction}\label{introduction}

Over the past decade, Unmanned Aerial Vehicles (UAVs) have rapidly
evolved from manually piloted platforms confined to open skies into
highly autonomous agents tasked with penetrating the most challenging,
inaccessible environments on Earth. Operational theaters such as
subterranean cave networks, structurally compromised industrial
facilities, dense urban canyons, and heavily foliated forests all share
a critical, defining vulnerability: the absolute denial, severe
degradation, or malicious spoofing of Global Navigation Satellite System
(GNSS) signals. Stripped of global positioning tethers, a micro-aerial
vehicle must rely entirely on its onboard suite of exteroceptive and
proprioceptive sensors to estimate its six-degree-of-freedom (6-DoF)
kinematic state, map unknown obstacles in real time, and execute
aggressive, collision-free maneuvers.

The engineering reality of achieving sustained autonomous flight in
these regimes is governed by brutal physical constraints. Micro- and
small-UAVs are bound by unforgiving Size, Weight, Power, and Cost
(SWaP-C) budgets. Every gram allocated to advanced compute modules
(e.g., NVIDIA Jetson architectures) or active perception sensors (e.g.,
mechanical LiDAR) directly penalizes flight endurance, which typically
hovers between 12 and 28 minutes. Simultaneously, the state estimation
algorithms running on these processors face unrelenting environmental
hostilities: dynamic illumination shifts that blind optical cameras,
featureless corridors that induce geometric degeneracy in LiDAR scans,
and rotor-induced downwash that violently agitates surrounding foliage.

While the academic literature addressing GPS-denied aerial navigation
has exploded in volume, the secondary literature attempting to organize
it remains fragmented. Existing survey articles are predominantly
qualitative, heavily biased toward specific sensor silos (such as pure
monocular vision), and critically lack transparent, reproducible data
provenance. To resolve these ambiguities and establish a definitive
baseline for the community, this systematic review conducts an
exhaustive, data-grounded synthesis of 279 manually extracted,
peer-reviewed publications spanning from the foundational efforts of
2013 through the swarm and neural paradigms of mid-2026.

\subsubsection{Research Questions}\label{research-questions}

This review targets four primary research questions designed to dissect
the architectural, environmental, algorithmic, and operational realities
of GPS-denied flight: * \textbf{RQ1 (Sensor Modalities \& Fusion
Trends):} Which sensor configurations dominate GPS-denied UAV
navigation, and how have multi-sensor fusion paradigms shifted
chronologically over the past decade? * \textbf{RQ2 (Localization
Accuracy \& Operational Environments):} What are the true quantitative
localization accuracy limits, drift rates, and robustness metrics
achieved across distinct degradation regimes, including indoor
warehouses, subterranean tunnels, and contested airspace? * \textbf{RQ3
(Algorithmic Paradigms \& Validation Fidelity):} Which algorithmic
frameworks lead the state of the art, and how does the maturity of their
physical flight validation compare against purely numerical simulation?
* \textbf{RQ4 (Systemic Limitations \& Future Agenda):} What are the
fundamental SWaP-C bottlenecks and perception vulnerabilities documented
across the literature, and what open research priorities must be
addressed to unlock certifiable real-world deployment?

\begin{center}\rule{0.5\linewidth}{0.5pt}\end{center}

\subsection{2. Related Work \& Systematic
Positioning}\label{related-work-systematic-positioning}

Prior secondary literature has illuminated individual facets of UAV
state estimation. However, a rigorous audit of the existing review
articles identified within our screened corpus highlights severe
methodological gaps:

Unlike preceding qualitative surveys, the present study asserts its
authority through four methodological pillars: 1. It adheres strictly to
the \textbf{PRISMA 2020} methodology, enforcing reproducible inclusion
boundaries across IEEE Xplore and Scopus databases. 2. It evaluates
every publication against an objective \textbf{Plan-Do-Check-Act (PDCA)
Quality Appraisal} model, exposing the underlying rigor, baseline
comparisons, and reproducibility of the field. 3. It completely forbids
the artificial statistical pooling of heterogeneous metrics, reporting
instead verifiable numerical performance ranges exactly as they appear
in the primary literature. 4. It synthesizes cross-cutting physical
trade-offs across a massive, manually verified 279-paper gold-standard
evidence base.

\begin{center}\rule{0.5\linewidth}{0.5pt}\end{center}

\subsection{3. Systematic Review
Methodology}\label{systematic-review-methodology}

\subsubsection{3.1 Literature Search Strategy \& Information
Sources}\label{literature-search-strategy-information-sources}

The literature search was executed on June 15, 2026, targeting the two
premier indexing databases for robotics and aerospace engineering:
\textbf{IEEE Xplore} and \textbf{Scopus}. The automated search retrieved
exactly 2,000 raw candidate records, ensuring an exhaustively broad
initial net.

\subsubsection{3.2 Screening Protocol \& Eligibility
Criteria}\label{screening-protocol-eligibility-criteria}

Following automated cross-database deduplication, candidate records were
screened against strict operational criteria (detailed in
\textbf{{[}Insert Table I: Inclusion/exclusion criteria{]}}).

\textbf{Table I: Inclusion and Exclusion Criteria} \textbar{} Dimension
\textbar{} Inclusion Criteria \textbar{} Exclusion Criteria \textbar{}
\textbar---\textbar---\textbar---\textbar{} \textbar{} \textbf{Date}
\textbar{} Published 2010-01-01 to 2026-06-15 \textbar{} Pre-2010
publications \textbar{} \textbar{} \textbf{Language} \textbar{} English
\textbar{} Non-English \textbar{} \textbar{} \textbf{Document Type}
\textbar{} Peer-reviewed journal or full conference \textbar{}
Preprints, patents, theses, tutorials \textbar{} \textbar{}
\textbf{Topic} \textbar{} GNSS-denied/degraded aerial navigation
\textbar{} GNSS-dependent systems \textbar{} \textbar{}
\textbf{Platform} \textbar{} Unmanned Aerial Vehicles (UAVs) \textbar{}
Terrestrial (UGV) or underwater platforms \textbar{} \textbar{}
\textbf{Method} \textbar{} Multi-sensor fusion (≥2 modalities)
\textbar{} Unassisted single-sensor odometry \textbar{} \textbar{}
\textbf{Validation} \textbar{} Physical flight or high-fidelity
simulation \textbar{} Purely conceptual without empirical data
\textbar{}

Through multi-stage title, abstract, and full-text screening, exactly
\textbf{279 studies} met all eligibility thresholds.

The complete flow of information through the different phases of a
systematic review is depicted in the PRISMA 2020 flow diagram
(\textbf{{[}Insert Fig. 1: PRISMA flow diagram{]}}).

\subsubsection{3.3 Data Extraction \& PDCA Quality
Appraisal}\label{data-extraction-pdca-quality-appraisal}

Data extraction captured operational environments, platform dynamics,
sensor configurations, and verifiable performance outcomes. The full
evidence matrix for all included studies is provided in
\textbf{{[}Insert Table II: Evidence matrix (all 279 papers){]}}.

Quality screening evaluated each study across Methodological Rigor,
Reporting Quality, Baseline Comparative Rigor, and Reproducibility.
Stratification yielded \textbf{3 Q-High studies (1.1\%)}, \textbf{98
Q-Medium studies (35.1\%)}, and \textbf{178 Q-Low studies (63.8\%)},
establishing an average corpus composite quality score of 3.81 out of
8.0.

\begin{center}\rule{0.5\linewidth}{0.5pt}\end{center}

\subsection{4. Systematic Results \& Empirical Evidence
Synthesis}\label{systematic-results-empirical-evidence-synthesis}

\subsubsection{4.1 Chronological \& Geographic
Trajectory}\label{chronological-geographic-trajectory}

The temporal trajectory of the corpus illustrates a steady expansion
that accelerated exponentially after 2021 (\textbf{{[}Insert Fig. 2:
Publications per year{]}}). Geographically, research output is heavily
concentrated across major robotics ecosystems in China, the United
States, and Singapore.

\subsubsection{4.2 RQ1 Synthesis: Sensor Modalities \& Fusion
Architectural
Shifts}\label{rq1-synthesis-sensor-modalities-fusion-architectural-shifts}

To answer \textbf{RQ1}, we analyzed the operational prevalence of
primary sensor modalities (\textbf{{[}Insert Fig. 3: Sensor
distribution{]}}) and fusion configurations across three historical
epochs.

\textbf{{[}Insert Table III: Method comparison summary (Evolution of
Fusion Architectures){]}} \textbf{Table III: Method Comparison Summary}
\textbar{} Multi-Sensor Pairing \textbar{} 2013--2016 (\%) \textbar{}
2017--2021 (\%) \textbar{} 2022--2026 (\%) \textbar{} Overall Frequency
\textbar{}
\textbar---\textbar---\textbar---\textbar---\textbar---\textbar{}
\textbar{} Camera + IMU \textbar{} 34.6\% \textbar{} 27.8\% \textbar{}
23.0\% \textbar{} 25.4\% \textbar{} \textbar{} Camera + LiDAR \textbar{}
15.4\% \textbar{} 16.5\% \textbar{} 17.2\% \textbar{} 16.8\% \textbar{}
\textbar{} LiDAR + IMU \textbar{} 15.4\% \textbar{} 8.9\% \textbar{}
17.2\% \textbar{} 14.7\% \textbar{} \textbar{} Camera + LiDAR + IMU
\textbar{} 7.7\% \textbar{} 5.1\% \textbar{} 10.3\% \textbar{} 8.6\%
\textbar{} \textbar{} UWB + IMU \textbar{} 0.0\% \textbar{} 7.6\%
\textbar{} 10.9\% \textbar{} 9.0\% \textbar{}

\textbf{Key Findings for RQ1:} 1. \textbf{The Obsolescence of Ultrasonic
Sensing vs.~Rise of UWB:} In the foundational era (2013--2016),
downward-facing sonars were prevalent (30.8\%) for altitude hold. By
2022--2026, ultrasonic sensors almost vanished (0.6\%), aggressively
displaced by lightweight Time-of-Flight (ToF) LiDAR and Ultra-Wideband
(UWB) networks. 2. \textbf{Solidification of Visual-Inertial \&
LiDAR-Inertial Couplings:} Camera + IMU remains the standard baseline
pairing. However, LiDAR + IMU and triple-sensor arrays (Camera + LiDAR +
IMU) grew five-fold in raw frequency after 2021, unlocked by the
miniaturization of solid-state LiDAR scanners operating within strict
micro-UAV payload allowances.

\subsubsection{4.3 RQ2 Synthesis: Environmental Degradation Regimes \&
Metric
Envelopes}\label{rq2-synthesis-environmental-degradation-regimes-metric-envelopes}

To address \textbf{RQ2}, the corpus was categorized into operational
environments (\textbf{{[}Insert Fig. 5: Environment distribution{]}}),
extracting verifiable trajectory error metrics.

\textbf{{[}Insert Table IV: Key results summary by environment{]}}
\textbf{Table IV: Key Results Summary by Environmental Degradation
Regime} \textbar{} Environment Category \textbar{} Primary Degradation
Stressor \textbar{} Reported Accuracy Envelopes \textbar{} Key Example
Reference \textbar{}
\textbar---\textbar---\textbar---\textbar---\textbar{} \textbar{} Mixed
/ Outdoor Denied \textbar{} GNSS dropouts, altitude variance \textbar{}
Horizontal RMSE: 0.10 m -- 1.84 m \textbar{} {[}REC\_0037{]} \textbar{}
\textbar{} Indoor \& Warehouses \textbar{} Geometric symmetry, multipath
\textbar{} Mean path error: 0.039 m -- 0.15 m \textbar{} {[}REC\_0028{]}
\textbar{} \textbar{} Urban Canyons \textbar{} Multipath, structural
shadow \textbar{} Mean position error: 0.80 m -- 2.0 m \textbar{}
{[}REC\_0312{]} \textbar{} \textbar{} Dense Forest Canopy \textbar{}
Dynamic foliage, non-rigid slip \textbar{} Obstacle tracking FPS: 30--40
Hz \textbar{} {[}REC\_1253{]} \textbar{} \textbar{} Subterranean/Tunnels
\textbar{} Zero illumination, airborne dust \textbar{} 3D error: 0.05 m
-- 2.6 m \textbar{} {[}REC\_1267{]} \textbar{}

Subterranean environments represent the most hostile operational regime,
characterized by total darkness, suspended dust, and featureless
cylindrical cross-sections that induce catastrophic geometric
degeneracy.

\subsubsection{4.4 RQ3 Synthesis: Algorithmic Paradigms \& Validation
Fidelity}\label{rq3-synthesis-algorithmic-paradigms-validation-fidelity}

To answer \textbf{RQ3}, we cross-tabulated algorithmic families
(\textbf{{[}Insert Fig. 4: Method category distribution{]}}) against
empirical validation fidelity (\textbf{{[}Insert Fig. 7: Real vs Sim
breakdown{]}}).

\textbf{Critical Algorithmic Insights:} 1. \textbf{The Swarm Sim-to-Real
Disconnect:} While cooperative swarm navigation represents 17.6\% of the
literature, it suffers from the lowest physical flight-testing rate in
the entire corpus: \textbf{49.0\% of swarm studies rely entirely on
numerical simulations}. Physical swarm validation remains stalled by
inter-agent RF packet loss and aerodynamic downwash collisions. 2.
\textbf{LiDAR SLAM Hardware Grounding:} Conversely, LiDAR-inertial SLAM
exhibits the highest hardware validation fidelity: \textbf{85.7\% of
LiDAR studies execute physical flight tests}, driven by mature, robust
open-source stacks.

\subsubsection{4.5 The Three Q-High Benchmark Anchor
Studies}\label{the-three-q-high-benchmark-anchor-studies}

Only 3 of the 279 studies (1.1\%) met all criteria for \textbf{Q-High}
classification, representing the absolute pinnacle of benchmarking,
rigor, and reproducibility in the field: 1. \textbf{REC\_0037 (2026):}
Extracted deep learned descriptors (SuperPoint/LightGlue) between
onboard monocular video and georeferenced satellite orthoimagery,
achieving a mean horizontal error of 1.84 m across multi-kilometer
flights. 2. \textbf{REC\_0502 (2026):} Deployed 3D Gaussian Splatting
combined with a 200 Hz IMU to achieve an unprecedented 0.80 m mean error
at high altitudes, vastly outperforming ORB-SLAM3 baselines. 3.
\textbf{REC\_1277 (2026):} Utilized physics-guided neural networks
(PG-TLNet) to compensate for aggressive aerodynamic and electromagnetic
interference across a fixed-wing swarm, reducing magnetic interference
standard deviation by 85.2\%.

\begin{center}\rule{0.5\linewidth}{0.5pt}\end{center}

\subsection{5. Discussion \& Cross-Cutting
Tensions}\label{discussion-cross-cutting-tensions}

\subsubsection{5.1 The Accuracy--SWaP--Compute
Trilemma}\label{the-accuracyswapcompute-trilemma}

The synthesis reveals that GPS-denied aerial navigation is governed by
an inescapable physical trilemma. High-accuracy systems employing
multi-beam LiDAR arrays achieve sub-decimeter accuracy but demand
massive power (20--65 W) and decimate flight endurance. Conversely,
ultra-lightweight nano-MAVs equipped with mere monocular cameras
accumulate unbounded dead-reckoning drift, causing catastrophic
localization divergence on extended flights. The intermediate
solution---deep neural matching and radiance field
rendering---eliminates drift but introduces severe compute latency
bottlenecks, restricting closed-loop flight speeds.

\subsubsection{5.2 Metric Heterogeneity \& Non-Standardized
Evaluation}\label{metric-heterogeneity-non-standardized-evaluation}

A severe systemic weakness identified in the corpus is the extreme
divergence in performance metric reporting. Over \textbf{90\% of the
corpus evaluates algorithms exclusively on private, self-collected
flight sequences}, severely undermining comparative analysis.
Furthermore, researchers arbitrarily alternate between reporting
Absolute Trajectory Error (ATE), Relative Pose Error (RPE), or
percentage drift, rendering direct algorithmic comparisons impossible
without standardized benchmarks.

\subsubsection{5.3 The Sim-to-Real Disconnect \& Offline
Trap}\label{the-sim-to-real-disconnect-offline-trap}

Over \textbf{51\% of all studies in the corpus never execute closed-loop
autonomous flight}. A widespread paradigm involves recording sensor
rosbags during manual radio-controlled flight, executing SLAM estimation
offline on desktop workstations, and reporting post-processed accuracy.
This offline methodology dangerously ignores real-world compute latency,
rotor vibration harmonics, and aerodynamic ground-effect disturbances.

\begin{center}\rule{0.5\linewidth}{0.5pt}\end{center}

\subsection{6. RQ4 Synthesis: Technological Limitations \& Future
Research
Agenda}\label{rq4-synthesis-technological-limitations-future-research-agenda}

To address \textbf{RQ4}, we systematically extracted and categorized the
failure modes, technological limitations, and proposed future directions
across all 279 studies.

The top documented limitations revolve around visual degradation in
extreme lighting, edge GPU thermal throttling, cumulative dead-reckoning
drift in feature-sparse environments, and RF packet loss in multi-agent
swarms.

\subsubsection{Strategic Future Research
Directions:}\label{strategic-future-research-directions}

\begin{enumerate}
\def\labelenumi{\arabic{enumi}.}
\tightlist
\item
  \textbf{Neuromorphic \& Bio-Inspired Event Sensing:} Integrating
  event-based neuromorphic cameras operating at microsecond temporal
  resolution offers a viable path toward blur-free state estimation
  under high-speed aggressive flight, defeating the limitations of
  standard CMOS sensors.
\item
  \textbf{Decentralized, Communication-Resilient Swarm SLAM:} Future
  swarm algorithms must implement communication-censored factor graphs,
  maintaining formation stability during extended RF dropouts by
  transmitting only high-information marginal keyframes.
\item
  \textbf{Physics-Informed Neural Network (PINN) Sensor Compensation:}
  Physics-guided neural models must be leveraged to predict and
  compensate for complex non-linear motor interference, thermal IMU bias
  drifts, and aerodynamic rotor downwash dynamics in real time.
\item
  \textbf{Community Standardization \& Open-Source Verification:} The
  robotics community must urgently establish standardized aerial
  navigation benchmarks providing synchronous hardware data coupled with
  sub-millimeter motion-capture ground truth in hostile subterranean and
  canopy environments.
\end{enumerate}

\begin{center}\rule{0.5\linewidth}{0.5pt}\end{center}

\subsection{7. Conclusions}\label{conclusions}

This systematic literature review provides the most comprehensive,
data-grounded synthesis of GPS-denied UAV navigation conducted to date,
analyzing 279 gold-standard peer-reviewed studies published between 2013
and mid-2026. By examining the corpus across four foundational research
questions, we have mapped the operational obsolescence of legacy
sensors, defined precise localization performance envelopes across
hostile environments, exposed a profound simulation dependency in
cooperative swarm literature, and articulated the systemic
vulnerabilities imposed by strict SWaP-C budgets.

By adhering strictly to transparent data extraction, avoiding artificial
metric pooling, and grounding all findings in verifiable evidence, this
review provides a definitive reference architecture for the next
generation of resilient, certifiably autonomous aerial robotics.

\begin{center}\rule{0.5\linewidth}{0.5pt}\end{center}

\subsection{References}\label{references}

\emph{(Full IEEE bibliographic citations mapped to all 279 in-corpus
records; primary anchors and representative studies referenced directly
via \texttt{07\_manuscript/references.bib}).}

\bibliographystyle{IEEEtran}
\bibliography{references}

\end{document}
